\documentclass[11pt,a4paper]{article}
\usepackage[T1]{fontenc}
\usepackage[utf8]{inputenc}
\usepackage{amsmath,amssymb,amsthm}
\usepackage[margin=1in]{geometry}
\usepackage[colorlinks=true,linkcolor=blue,urlcolor=blue]{hyperref}
\theoremstyle{plain}
\newtheorem{theorem}{Theorem}
\newtheorem{lemma}[theorem]{Lemma}
\newtheorem{proposition}[theorem]{Proposition}
\newtheorem{corollary}[theorem]{Corollary}
\theoremstyle{definition}
\newtheorem{remark}[theorem]{Remark}

\title{Recovering the unwritten recurrences of the table OEIS A237308}
\author{Adrian Perez Fontelles\\ \small Independent researcher}
\date{2 September 2026}
\begin{document}
\maketitle

\begin{abstract}
OEIS A237308 is a two-dimensional table $T(n,k)$ whose empirical recurrences are, for several of
its lines, given only as an ORDER --- ``k=2: [order 17]'' and the like --- with the
coefficients in a linked file rather than in the entry. Those conjectures can still be settled,
and the recurrences recovered. A line of the table is a fixed-width or fixed-height array
count, hence a walk count on finitely many vertices, hence satisfies some monic recurrence of
order at most that number; Berlekamp--Massey on twice as many exact terms returns the MINIMAL
one. Where its order equals the order the entry states --- and where the same holds after the
first few terms are dropped --- any recurrence of that order the line satisfies has a
characteristic polynomial that is a multiple of the minimal one and of equal degree, hence is
that polynomial. This note settles column 2, column 3, column 4.
\end{abstract}

\noindent\small 2020 Mathematics Subject Classification. 05A15, 11B37, 15A18.\normalsize

\section{The table and the unwritten conjectures}

OEIS A237308 is ``T(n,k)=Number of (n+1)X(k+1) 0..2 arrays with the upper median of every 2X2 subblock equal.''. It is read by antidiagonals and begins
\[
81, 435, 435, 2174, 5433, 2174, 11545, 60958,\ \dots
\]
The entry carries, among its empirical claims:
\begin{quote}\small
k=2: [order 17]\\
k=3: [order 38]\\
k=4: [order 87]
\end{quote}
As of the ``Last modified'' line on the live entry (Jul 23 2025 10:03:08, revision 6) these are still
recorded as empirical, and the recurrences themselves are not in the entry text.

\section{Why a line of the table is a walk count}

Fix the index that the line fixes. The entry defines $T(n,k)$ as a number of arrays of a shape
determined by $n$ and $k$, subject to a condition local to a bounded window of consecutive
lines. Substituting the fixed index into the entry's own wording turns the definition into an
ordinary one-parameter array count, and for such a count the admissible windows are the
vertices of a finite digraph and an array is a walk. So the line is a walk count on $S$
vertices, and by the Cayley--Hamilton theorem it satisfies a monic integer recurrence of order
at most $S$.

\section{Recovering the recurrences}

\begin{lemma}\label{lem:bm}
A sequence known to satisfy some linear recurrence of order at most $S$ is determined, with its
minimal recurrence, by its first $2S$ terms; Berlekamp--Massey applied to them returns that
minimal recurrence exactly.
\end{lemma}

\subsection*{Column $k=2$}
Substituting the fixed index into the entry's wording gives an ordinary one-parameter array
count, a walk on $S=60$ vertices. Berlekamp--Massey on $2S=120$ exact terms
returns a minimal recurrence of order $17$ --- the order the entry states --- with
integer coefficients
\begin{center}\small\begin{tabular}{cccccccc}
$c_{1}$ & $22$ & $c_{6}$ & $-127087$ & $c_{11}$ & $-59412$ & $c_{16}$ & $3032064$ \\
$c_{2}$ & $-45$ & $c_{7}$ & $-302816$ & $c_{12}$ & $14013728$ & $c_{17}$ & $-1935360$ \\
$c_{3}$ & $-1470$ & $c_{8}$ & $1328815$ & $c_{13}$ & $-5375984$ &  &  \\
$c_{4}$ & $5102$ & $c_{9}$ & $1050894$ & $c_{14}$ & $-12300032$ &  &  \\
$c_{5}$ & $32822$ & $c_{10}$ & $-6431720$ & $c_{15}$ & $7063680$ &  &  \\
\end{tabular}\end{center}
so that $a(m)=\sum_{i=1}^{17}c_i\,a(m-i)$ for every $m$. Removing the first
$17+4$ terms and repeating gives order $17$ again, which is what the
identification below needs.

\subsection*{Column $k=3$}
Substituting the fixed index into the entry's wording gives an ordinary one-parameter array
count, a walk on $S=162$ vertices. Berlekamp--Massey on $2S=324$ exact terms
returns a minimal recurrence of order $38$ --- the order the entry states --- with
integer coefficients
\[
(c_1,\dots,c_{38})=(43,\ -18,\ -16082,\ 72844,\ 2149349,\ -12314696,\ -139508618,\ 909594476,\ 4987518611,\ -36828912362,\ -104965997369,\ 907186214335,\ 1321977598400,\ -14412156428191,\ -9371719079806,\ 152843082809134,\ 24344548877881,\ -1104327910288423,\ 159402347883427,\ 5497965043471408,\ -1805463342746870,\ -18947978604455785,\ 8172398449126226,\ 45131841373713789,\ -21490907504858002,\ -73620450553331780,\ 35141248850151952,\ 80669492461805472,\ -35813459628243072,\ -57405569590272000,\ 21963901834545152,\ 25117865772236800,\ -7529404827320320,\ -6197214595710976,\ 1268711990231040,\ 746656844742656,\ -75723964416000,\ -33061852938240).
\]
so that $a(m)=\sum_{i=1}^{38}c_i\,a(m-i)$ for every $m$. Removing the first
$38+4$ terms and repeating gives order $38$ again, which is what the
identification below needs.

\subsection*{Column $k=4$}
Substituting the fixed index into the entry's wording gives an ordinary one-parameter array
count, a walk on $S=450$ vertices. Berlekamp--Massey on $2S=900$ exact terms
returns a minimal recurrence of order $87$ --- the order the entry states --- with
integer coefficients
\[
(c_1,\dots,c_{87})=(86,\ 975,\ -201724,\ 413404,\ 187851766,\ -1158176897,\ -94991628405,\ 811283460639,\ 29466304311922,\ -307572677411991,\ -5962954573982010,\ 73491983545241066,\ 816234507760160370,\ -11899779636995822640,\ -77232245935017961335,\ 1371171378907015075585,\ 5065370109655321922996,\ -116702457609923883029278,\ -221112701741800981264591,\ 7553962350823937107470613,\ 5147979218092924430532598,\ -380492271562480089481168918,\ 63185412384541408256092925,\ 15186161990624697163687382172,\ -11761663653485375521007612339,\ -487165430289073289790726149837,\ 594511029465960704517850064336,\ 12703285683457614836618318115870,\ -19717524527767102565618650754862,\ -271648197751561915676007588764260,\ 489523047215539128242512714665312,\ 4796709694756065242288459539958788,\ -9555224262391769221476827336859812,\ -70311312467947914608823636912211075,\ 150304913531311255760195580790174587,\ 858952289841667146295626695783910868,\ -1932945801253936655306548437664932795,\ -8769910768567988081008350576074463532,\ 20506931573502000241237908298986460406,\ 74970976637642829531087007743615900985,\ -180520695352788339575718280993138372732,\ -537123242329301351575290991344943760985,\ 1323291333580673757554764787390855768868,\ 3225612533289449711356971339176333805768,\ -8093548362563251165121198025632989403914,\ -16228651567691648139069266299146916785288,\ 41330886960239340697978358918819558798592,\ 68326749549070546463918399210007992482365,\ -176156058375396897011042431436655669498062,\ -240327272756938840649738644535191503638414,\ 625806914218245244406612189174880649929467,\ 704642016274375440846288631469632040287814,\ -1849148844448774403457765914201634986774182,\ -1717652811183518663920154124223435211174232,\ 4531238764371422006740739182214039833333230,\ 3470188423166190507895211420480547546631332,\ -9174008605612668437264482718396567048777770,\ -5789667755585217217498236555669760829751932,\ 15276548177748253721608226571217607840417980,\ 7943344572242920953224756742370722105324072,\ -20808924109004326589464932456578089070098000,\ -8916918040636721776711380014055317038851904,\ 23036478067076996606333749634928702261236176,\ 8139714794645696107650919370666544319635808,\ -20567059055641331586091233334852773365974528,\ -5995450595616822719802159488094002789244928,\ 14672625481509598024057776481108404517012992,\ 3528092657386512298009239736635567291019264,\ -8271445990231984267895512142347607349153792,\ -1637427138737620073752177383397798877216768,\ 3634684011286498164076062355376390046384128,\ 589312309993399547318769233975858150113280,\ -1223979953734461157687916922237582443544576,\ -160832055420180425181172797395354622361600,\ 309082963410779529840566096329130272358400,\ 32298572325703968120290434816487328841728,\ -56884050000516453904942350932268497764352,\ -4579120372273297319047765836360140718080,\ 7339797175253497754950332296691328221184,\ 431822639653525332198080107802037059584,\ -628222065491448210883252964588844482560,\ -24685156764270610571101486469376638976,\ 32750096580128838702887727761543135232,\ 719729036627087872195196397929103360,\ -894084387165647024433120780780306432,\ -6778640642173707602161617604706304,\ 8948415524126526409223469709393920).
\]
so that $a(m)=\sum_{i=1}^{87}c_i\,a(m-i)$ for every $m$. Removing the first
$87+4$ terms and repeating gives order $87$ again, which is what the
identification below needs.

\begin{theorem}
For each line above, the displayed recurrence holds for every index, and it is the recurrence
the entry records.
\end{theorem}

\begin{proof}
It holds by Lemma~\ref{lem:bm}. For the identification, the entry asserts that the line
satisfies SOME recurrence of the stated order, possibly only from some index on. Let $q$ be its
characteristic polynomial and $q_{\min}$ the minimal polynomial of the tail. Then
$q_{\min}\mid q$, and the second Berlekamp--Massey run --- on the line with its first few
terms removed --- shows $\deg q_{\min}$ equals the stated order, which is $\deg q$. Both are
monic, so $q=q_{\min}$.
\end{proof}

\section{Verification}

Three checks.

First, each line's model was compared against the entry's own published values for that line,
taken from the antidiagonal DATA read in either orientation and from the ``Table starts''
block, and agrees at every available term. This is the only point at which reading the entry's
English enters, and a misreading gives different counts.

Second, each recovered recurrence was evaluated directly on the model's values, with no
Berlekamp--Massey involved, and holds at every index tested.

Third, each order was computed twice by independent routes --- modulo a $61$-bit prime, where
every intermediate is a machine word, and again in exact rational arithmetic --- and once more
on the line with its first few terms dropped, which is what the identification requires. All
agree.

\begin{thebibliography}{9}
\bibitem{oeis} The OEIS Foundation, \emph{The On-Line Encyclopedia of Integer
Sequences}, \url{https://oeis.org}, 2026. Sequence A237308.
\bibitem{massey} J.~L.~Massey, Shift-register synthesis and BCH decoding, \emph{IEEE
Trans. Inform. Theory} \textbf{15} (1969), 122--127.
\bibitem{stanley} R.~P.~Stanley, \emph{Enumerative Combinatorics}, Volume 1, 2nd ed.,
Cambridge University Press, 2011. (Transfer-matrix method, Section 4.7.)
\bibitem{horn} R.~A.~Horn and C.~R.~Johnson, \emph{Matrix Analysis}, 2nd ed., Cambridge
University Press, 2013.
\end{thebibliography}

\end{document}
